\chapter{Sleep and Wakeup}
\label{CH:SLEEP}
%%

Scheduling and locks help conceal the actions of one thread
from another,
but we also need abstractions that help
threads intentionally interact.
For example, the reader of a pipe in xv6++ may need to wait
for a writing process to produce data;
a parent's call to \lstinline{wait} may need to
wait for a child to exit; and
a process reading the disk needs to wait
for the disk hardware to finish the read.
The xv6++ kernel uses a mechanism called sleep and wakeup
in these situations (and many others).
Sleep allows a kernel thread to
wait for some condition to be true; another thread
or an interrupt handler can cause the condition
to be true (typically by modifying some variable(s))
and then call wakeup
to indicate that threads waiting for the condition should resume.
Sleep and wakeup are often called
\indextext{sequence coordination}
or
\indextext{conditional synchronization}
mechanisms.

Before proceeding, please read \texttt{kernel/process\_manager.h},
\texttt{sleep} and \texttt{wakeup}, and the pipe implementation in
\texttt{kernel/pipe.cpp}.

\section{Overview}

The sleep/wakeup interface looks like:

\begin{lstlisting}
void process::sleep(void *chan, spin_lock &lk)
void process_manager::wakeup(void *chan)
\end{lstlisting}

\lstinline{sleep()} marks the calling process as \lstinline{SLEEPING}
(not \lstinline{RUNNABLE}) and releases the CPU by context-switching
to the scheduler, so that other processes can run.
The \lstinline{chan} argument is called the \indextext{wait channel}.
\lstinline{wakeup(chan)} wakes up all processes (if any)
that have called \lstinline{sleep(chan, ...)} with the same \lstinline{chan}
value. \lstinline{sleep} and \lstinline{wakeup} treat \lstinline{chan}
as an opaque 64-bit value; the only thing they do with it is compare for
equality. The usual pattern is for callers to pass the address of
some convenient object as the \lstinline{chan} argument.

Kernel code calls \lstinline{sleep} to wait for some
\indextext{condition} to become true. For example, the kernel code
that reads from a pipe calls \lstinline{sleep} if the pipe buffer is
currently empty; the condition in this case is the pipe buffer
becoming non-empty (due to another process writing to the pipe).
\lstinline{sleep} and \lstinline{wakeup} do not know what the
condition is: only the calling code knows. The usual pattern is for
the caller to first check the condition, and call \lstinline{sleep}
if it is not true; code that later makes the
condition true calls \lstinline{wakeup}.

Here's a sketch of how the xv6++ pipe code uses
\lstinline{sleep} and \lstinline{wakeup}:

\begin{lstlisting}
int pipe::read(uint64 addr, int n) {
  auto pr = kernel.cpus.curproc();
  lock_guard<spin_lock> g(lock);

  while (nread == nwrite && writeopen) {
    (*@\textcolor{blue}{// ZZZ}@*)
    pr->sleep(&nread, lock);
  }

  remove data from data[nread % PIPESIZE];
  kernel.processes.wakeup(&nwrite);
}

int pipe::write(uint64 addr, int n) {
  auto pr = kernel.cpus.curproc();
  lock_guard<spin_lock> g(lock);

  while (nwrite == nread + PIPESIZE) {
    kernel.processes.wakeup(&nread);
    pr->sleep(&nwrite, lock);
  }

  append data to data[nwrite % PIPESIZE];
  kernel.processes.wakeup(&nread);
}
\end{lstlisting}

This code uses the address of
\lstinline{nread}
as the reader wait channel and the address of
\lstinline{nwrite}
as the writer wait channel.

What is the \lstinline{lk} argument to \lstinline{sleep}?
In all uses of \lstinline{sleep}/\lstinline{wakeup}
the condition involves shared data, used by both the
thread that sleeps and the thread that calls \lstinline{wakeup}, so
there always turns out to be a lock that protects the condition. That lock is
called the \indextext{condition lock}. In the pipe code above, both
functions use the pipe and its buffer while holding the pipe lock,
which in this case is also the condition lock. It's a rule that any code
that calls \lstinline{sleep} or \lstinline{wakeup} must hold the
condition lock, and that the lock must be passed to \lstinline{sleep}
as the second argument.

The reason that the condition lock must be held when \lstinline{sleep}
is called, and that it must be passed to \lstinline{sleep}, is to
prevent the possibility that another thread might call
\lstinline{wakeup} between the check of the condition
and the call to \lstinline{sleep}. A call to
\lstinline{wakeup} at that point would find no sleeping process to
wake up; the \lstinline{wakeup} would simply return.
But then the call to \lstinline{sleep} might never wake up,
since the \lstinline{wakeup} intended for it has already happened.
This undesirable situation is called a \indextext{lost wake-up}.

In the pipe example above, the lost wake-up being avoided is the
possibility that a thread on another CPU might call \lstinline{pipe::write}
at the point marked {\tt ZZZ},
between \lstinline{pipe::read}'s check of the condition and its call to
\lstinline{sleep}. The fact that \lstinline{pipe::read} holds the pipe
lock during the time between when it checks the condition and calls
\lstinline{sleep} prevents \lstinline{pipe::write} from executing, and
thus prevents a lost wake-up.

\lstinline{sleep()} releases the condition lock so that the code calling
\lstinline{wakeup()} can proceed. \lstinline{sleep()} also context-switches to the
scheduler in order to let other threads run while it is waiting. The
implementation performs these two steps in a way that is atomic
(indivisible) with respect to \lstinline{wakeup()}, to prevent lost
wake-ups.

%%
\section{Code: Sleep and wakeup}
%%

xv6++ implements the interface used in the example above with
\indexcode{sleep}
and
\indexcode{wakeup}.
The code is in
\texttt{kernel/process\_manager.cpp}.
The basic idea is to have
\lstinline{sleep}
mark the current process as
\indexcode{SLEEPING}
and then call
\indexcode{sched}
to release the CPU;
\lstinline{wakeup}
looks for a process sleeping on the given wait channel
and marks it as
\indexcode{RUNNABLE}.

\lstinline{sleep}
acquires
\indexcode{p->lock}
and \textit{only then} releases
the condition lock
\lstinline{lk}.
The fact that \lstinline{sleep}
holds one or the other of these locks at all times is what
prevents a concurrent \lstinline{wakeup}
(which must acquire and hold both) from acting,
and thus prevents a lost wake-up.
Now that
\lstinline{sleep}
holds just
\lstinline{p->lock},
it can put the process to sleep by recording
the wait channel,
changing the process state to \texttt{SLEEPING},
and calling
\lstinline{sched}.
In a moment it will be clear why it's critical that
\lstinline{p->lock} is not released (by the scheduler) until after
the process is marked \texttt{SLEEPING}.

At some point, a process will acquire the condition lock,
set the condition that the sleeper is waiting for,
and call \lstinline{wakeup(chan)}.
It's important that \lstinline{wakeup} is called
while holding the condition lock\footnote{%
Strictly speaking it is sufficient if
\lstinline{wakeup}
merely follows the
\lstinline{acquire}
(that is, one could call
\lstinline{wakeup}
after the
\lstinline{release}).%
}.
\lstinline{wakeup}
loops over the process table.
It acquires the
\lstinline{p->lock}
of each process it inspects.
When \lstinline{wakeup} finds a process in state
\indexcode{SLEEPING}
with a matching
\indexcode{chan},
it changes that process's state to
\indexcode{RUNNABLE}.
The next time the scheduler runs, it will
see that the process is ready to be run.

\begin{figure}[t]
  \begin{tikzpicture}

  \node (code0) at (0,0) {
    \begin{lstlisting}[]
      pipe::read() {
        lock_guard<spin_lock> g(lock);
        while(nread == nwrite) {
          pr->sleep(&nread, lock) {
            // in process::sleep()
            lock_guard<spin_lock> pg(p->lock)
            lock.release()
            p->state = SLEEPING
            ...
            swtch() {
              // in scheduler()
              p->lock.release()
              ...
    \end{lstlisting}

    %% this line is important
   };

    \draw[<->,thick] (4.7,2.5) -- (4.7,-0.3);
    \draw[<->,thick] (4.9,0.7) -- (4.9,-2.6);
    \node at (6.5, 1.5) {Holding pipe lock};
    \node at (6.5, -1.3) {Holding p->lock};
\end{tikzpicture}

  \caption{Overlapping locks to avoid lost wake-up}
  \label{fig:overlap}
\end{figure}

Why do the locking rules for
\lstinline{sleep}
and
\lstinline{wakeup}
ensure that a process that's going to sleep
won't miss a concurrent wakeup?
The going-to-sleep process holds either
the condition lock or its own
\lstinline{p->lock}
or both from \textit{before} it checks the condition
until \textit{after} it has marked itself as \texttt{SLEEPING};
see Figure~\ref{fig:overlap}.
The process calling \texttt{wakeup} needs to acquire \textit{both}
locks.
The waker might acquire the locks first, which means
it will make the condition true before the consuming
thread checks the condition, and the consuming thread
won't need to call \lstinline{sleep()};
or the waker's \lstinline{acquire()}s might have to wait until the consuming
thread has completely finished going to sleep and
releases the locks, in which case the waker will then see
that the consuming thread is marked \texttt{SLEEPING}
and will wake it up.

Sometimes multiple processes are sleeping
on the same channel; for example, more than one process
reading from a pipe.
A single call to
\lstinline{wakeup}
will wake them all up.
One of them will run first and acquire the lock that
\lstinline{sleep}
was called with, and (in the case of pipes) read whatever
data is waiting.
The other processes will find that, despite being woken up,
there is no data to be read.
From their point of view the wakeup was ``spurious,'' and
they must sleep again.
For this reason \lstinline{sleep} is always called inside a loop that
re-checks the condition, as in \lstinline{pipe::read} above.

No harm is done if two uses of sleep/wakeup accidentally
choose the same channel: they will see spurious wakeups,
but looping as described above will tolerate this problem.
Much of the charm of sleep/wakeup is that it is both
lightweight (no need to create special data
structures to act as wait channels) and provides a layer
of indirection (callers need not know which specific process
they are interacting with).

%%
\section{Code: Pipes}
%%

xv6++'s pipes are an example of code that uses
\lstinline{sleep}
and \lstinline{wakeup} to synchronize producers and
consumers.
We saw the interface for pipes in Chapter~\ref{CH:UNIX}:
bytes written to one end of a pipe are copied
to an in-kernel buffer and then can be read from
the other end of the pipe.

Each pipe
is represented by a
\indexcode{class pipe},
which contains
a
\lstinline{lock}
and a
\lstinline{data}
buffer.
The fields
\lstinline{nread}
and
\lstinline{nwrite}
count the total number of bytes read from
and written to the buffer.
The buffer wraps around:
the next byte written after
\lstinline{data[PIPESIZE-1]}
is
\lstinline{data[0]}.
The counts do not wrap.
This convention lets the implementation
distinguish a full buffer from an empty buffer:
\begin{lstlisting}
nwrite == nread + PIPESIZE  // full
nwrite == nread             // empty
\end{lstlisting}
but it means that indexing into the buffer
must use
\lstinline{data[nread % PIPESIZE]}
instead of just
\lstinline{data[nread]}
(and similarly for
\lstinline{nwrite}).

Now suppose calls to
\lstinline{pipe::read}
and
\lstinline{pipe::write}
happen simultaneously on two different CPUs.
\lstinline{pipe::write}
begins by constructing a
\lstinline{lock_guard}
that acquires the pipe's lock, which
protects the counts, the data, and their
associated invariants.
\lstinline{pipe::read}
then tries to acquire the lock too, but cannot,
and spins waiting for the lock.

While
\lstinline{pipe::read}
waits,
\lstinline{pipe::write}
loops over the bytes being written,
adding each to the pipe in turn.
During this loop, it could happen that
the buffer fills.
In this case,
\lstinline{pipe::write}
calls
\lstinline{wakeup}
to alert any sleeping readers that there is data waiting,
and then sleeps (releasing the pipe lock) to wait for a reader
to take some bytes out of the buffer.

\lstinline{pipe::read}
now acquires the pipe lock and enters its critical section:
it finds that
\lstinline{nread != nwrite},
copies data out of the pipe, and increments
\lstinline{nread}
by the number of bytes copied.
That much space in the buffer is now available for writing, so
\lstinline{pipe::read}
calls
\lstinline{wakeup}
to wake any sleeping writers before it returns.

The pipe code uses separate wait channels for reader and writer
(\lstinline{&nread}
and
\lstinline{&nwrite}),
which can reduce spurious wakeups if there are multiple readers
and writers waiting for the same pipe.

%%
\section{Code: Wait, exit, and kill}
%%

Please read the code for the kernel functions
\texttt{process\_manager::wait}, \texttt{process::exit}, and
\texttt{process\_manager::kill}
in the process management code
(see \texttt{kernel/process\_manager.cpp}, \texttt{wait}),
(see \texttt{kernel/process.cpp}, \texttt{exit}),
and
(see \texttt{kernel/process\_manager.cpp}, \texttt{kill}).

\lstinline{sleep}
and
\lstinline{wakeup}
can be used for many kinds of waiting.
An interesting example is the interaction between a child's
\indexcode{exit}
and its parent's
\indexcode{wait}.
At the time of the child's death, the parent may already
be sleeping in \lstinline{wait}, or may be doing something else;
in the latter case, a subsequent call to \lstinline{wait} must
observe the child's death, perhaps long after it calls \lstinline{exit}.

xv6++ records a child's demise until \lstinline{wait}
observes it by putting the child in the \indexcode{ZOMBIE}
state, where it stays until the parent notices it, frees the child's
resources, and returns the child's process ID to the parent.
If the parent exits before the child, the parent gives the child to
\lstinline{init}, which perpetually calls \lstinline{wait};
thus every child has a parent to clean up after it.

\lstinline{process_manager::wait} starts by acquiring
\lstinline{wait\_lock},
which acts as the condition lock that helps ensure that
\lstinline{wait} doesn't miss a wakeup from an exiting child.
Then \lstinline{wait} scans the process table.
If it finds a child in \texttt{ZOMBIE} state,
it frees that child's resources, copies the child's exit status,
and returns the child's process ID.
If \lstinline{wait} finds children but none have exited,
it calls \lstinline{sleep} to wait for any of them to exit,
and then scans again.
\lstinline{wait} often holds two locks,
\lstinline{wait\_lock} and some process's \lstinline{p->lock};
the deadlock-avoiding order is first \lstinline{wait\_lock}
and then \lstinline{p->lock}.

\lstinline{process::exit} records the exit status, frees some resources,
calls a reparenting routine to give any children to \lstinline{init},
wakes up the parent in case it is in \lstinline{wait},
marks the caller as a zombie, and permanently yields the CPU.
\lstinline{process::exit} holds both \lstinline{wait\_lock} and \lstinline{p->lock}
during this sequence, in the same order as \lstinline{wait}, to avoid deadlock.

While \lstinline{exit} allows a process to terminate itself,
the \lstinline{kill} system call lets one process request that another terminate.
It would be too complex for \lstinline{kill} to directly destroy the victim
process, since the victim might be executing on another CPU.
Thus \lstinline{process_manager::kill} does very little: it sets the victim's killed flag
and, if it is sleeping, wakes it up.
Eventually the victim will enter or leave the kernel, at which point
trap handling code will notice that it has been killed and exit.

If the victim process is in \lstinline{sleep},
\lstinline{kill}'s call to \lstinline{wakeup} will cause the victim to return
from \lstinline{sleep}. This is potentially dangerous because the condition
being waited for may not be true. However, calls to \lstinline{sleep} are
always wrapped in a loop that re-tests the condition after \lstinline{sleep}
returns, and some loops also check the killed flag and abandon the operation
when that is correct.

%%
\section{Process locking}
%%

The lock associated with each process (\lstinline{p->lock}) is one of the
most important locks in xv6++.

A simple way to think about \lstinline{p->lock} is
that it must be held while reading or writing fields such as:
\lstinline{p->state},
\lstinline{p->chan},
\lstinline{p->killed},
\lstinline{p->xstate},
and
\lstinline{p->pid}.

However, most uses of \lstinline{p->lock} protect higher-level
invariants of xv6++'s process algorithms. For example:

\begin{itemize}

\item It prevents races while allocating process slots for new processes.

\item It conceals a process from view while it is being created
or destroyed.

\item It prevents a parent's \lstinline{wait} from collecting a process that
has set its state to \lstinline{ZOMBIE} but has not yet yielded the CPU.

\item It prevents another CPU's scheduler from deciding to run
a yielding process after it sets its state to \lstinline{RUNNABLE} but
before it finishes \lstinline{swtch}.

\item It ensures that only one CPU's scheduler decides to run a given
\lstinline{RUNNABLE} process.

\item Along with the condition lock, it helps prevent \lstinline{wakeup}
from overlooking a process that is calling \lstinline{sleep} but has not
finished yielding the CPU.

\item It prevents a \lstinline{kill}-target process from exiting and being
re-allocated between a \lstinline{kill} caller's checks and updates.

\end{itemize}

The \lstinline{p->parent} field is protected by the global lock
\lstinline{wait\_lock} rather than by \lstinline{p->lock}.
Only a process's parent modifies \lstinline{p->parent}, though
the field is read by the process itself and by other
processes searching for their children.
The purpose of \lstinline{wait\_lock} is to act as the condition lock when
\lstinline{wait} sleeps waiting for any child to exit.

%%
\section{Real world}
%%

\lstinline{sleep}
and
\lstinline{wakeup}
are a simple and effective synchronization method,
but there are many others;
semaphores~\cite{dijkstra65} are an example.
The first challenge in all of them is to
avoid the lost wakeups problem.

Because xv6++ runs on multiprocessors,
it uses explicit locks to make sleeping and waking atomic
with respect to each other.
Many real kernels use explicit per-channel or per-object wait queues
instead of scanning the whole process table in \lstinline{wakeup}.

xv6++'s \lstinline{wakeup} wakes all processes sleeping on a channel.
If there are many of them, most will discover that the condition is
no longer true and go back to sleep, wasting CPU time.
Thus many condition-variable designs provide two primitives:
\lstinline{signal} to wake one waiter, and \lstinline{broadcast} to wake all.

Forcibly killing processes also poses problems.
A killed process may be deep inside the kernel, and safely unwinding
requires care. Unix systems use signals, interrupted system calls,
and careful conventions to manage this complexity. xv6++ keeps things
simpler and thus inherits some limitations.
